\section{Estimación de máxima verosimilitud y otros métodos}
\begin{ejercicio}
    Sea $X\rightsquigarrow P_\theta$ con $\theta\in \mathbb{R}$ siendo $P_\theta$ una distribución con función de densidad
    \begin{equation*}
        f_\theta(x) = e^{\theta-x}, \qquad x\geq \theta
    \end{equation*}
    Dada una muestra aleatoria simple de tamaño $n$, encontrar los estimadores máximo verosímiles de $\theta$ y de $e^\theta$. Basándose en los resultados del Ejercicio~\ref{ej:7_rel4}, decir si estos estimadores son insesgados.\\

    \noindent
    Sea $(X_1, \ldots, X_n)$ una m.a.s. de $X\rightsquigarrow  P_\theta$, tenemos que:
    \begin{equation*}
        f^n_\theta(x_1, \ldots, x_n) \stackrel{\text{iid.}}{=} \prod_{i=1}^{n}f_\theta(x_i) = \prod_{i=1}^{n} e^{\theta-x} = e^{n\theta - \sum\limits_{i=1}^n x_i} \qquad \forall x_i \geq \theta
    \end{equation*}
    Por lo que si consideramos $x_{(1)} = \min\limits_{i \in \{1,\ldots,n\}}\{x_i\}$:
    \begin{equation*}
        L_{x_1,\ldots, x_n}(\theta) = f^n_\theta(x_1,\ldots,x_n) = \left\{\begin{array}{ll}
                e^{n\theta - \sum\limits_{i=1}^n x_i}  & \text{si\ } \theta \leq x_{(1)}  \\
             0 & \text{en otro caso\ } 
        \end{array}\right. 
    \end{equation*}
    Observamos que $L_{x_1,\ldots,x_n}(\theta)$ es una función creciente en $\left]-\infty,x_{(1)}\right]$, por lo que alcanza su máximo en $x_{(1)}$. Como el espacio paramétrico es $\Theta = \mathbb{R}$, tenemos que $\widehat{\theta}(X_1, \ldots, X_n) = X_{(1)}$ es un estimador de $\theta$, con lo que es el EMV, al maximizar la función de verosimilitud. Además, el Teorema de Zehna nos dice que:
    \begin{equation*}
        \widehat{e^\theta} = e^{\widehat{\theta}} = e^{X_{(1)}}
    \end{equation*}
    Por lo que $e^{X_{(1)}}$ es EMV para $e^{\theta}$.\\ 

    \noindent
    Buscamos ahora razonar si los estimadores obtenidos son o no insesgados, a partir de los resultados del Ejercicio~\ref{ej:7_rel4}.
    \begin{description}
        \item[Opción 1] Si suponemos que los EMVs que hemos obtenido son insesgados, vimos anteriormente en dicho ejercicio que el mínimo era suficiente y completo, por lo que tendríamos una función del suficiente y completo que es insesgada. Si comprobamos que también tiene momento de segundo orden finito y que es estimador (se hizo en el ejercicio anterior), tendríamos que es un UMVUE distinto al anterior, pero la unicidad del UMVUE nos dice que estos han de ser iguales, contradicción que viene de suponer que los EMVs son insesgados.
        \item[Opción 2] En el Ejercicio~\ref{ej:7_rel4} vimos que:
        \begin{itemize}
            \item El UMVUE de $\theta$ era:
            \begin{equation*}
                X_{(1)} - \frac{1}{n}
            \end{equation*}

            Como el UMVUE es insesgado, tenemos que:
            \begin{equation*}
                E\left[X_{(1)} - \frac{1}{n}\right] = \theta \Longrightarrow E[X_{(1)}] = \theta + \frac{1}{n} \neq \theta
            \end{equation*}

            Por lo que $X_{(1)}$ no es insesgado, y por lo tanto el EMV de $\theta$ no es insesgado.

            \item En dicho ejercicio también vimos que el UMVUE de $e^\theta$ era:
            \begin{equation*}
                \dfrac{n-1}{n}e^{X_{(1)}}
            \end{equation*}

            Como el UMVUE es insesgado, tenemos que:
            \begin{equation*}
                E\left[\dfrac{n-1}{n}e^{X_{(1)}}\right] = e^\theta \Longrightarrow E[e^{X_{(1)}}] = \frac{n}{n-1}e^\theta \neq e^\theta
            \end{equation*}

            Por lo que $e^{X_{(1)}}$ no es insesgado, y por lo tanto el EMV de $e^\theta$ no es insesgado.
        \end{itemize}
    \end{description}
\end{ejercicio}

\begin{ejercicio}
    Sea $(X_1, \ldots, X_n)$ una muestra aleatoria simple de una variable aleatoria $X$ con distribución exponencial. Basándose en los resultados del Ejercicio~\ref{ej:9_rel4}, encontrar los estimadores máximo verosímiles de la media y de la varianza de $X$.\\

    \noindent
    Sea $X\rightsquigarrow \exp(\lm)$ con $\lm\in \mathbb{R}^+$, nos piden encontrar EMVs para:
    \begin{equation*}
        E[X] = \frac{1}{\lm}, \qquad Var[X] = \frac{1}{\lm^2}
    \end{equation*}
    En el Ejercicio~\ref{ej:9_rel4} obtuvimos que dicha familia de distribuciones es regular, con $0<I_X(\lm)< \infty$. Además, vimos que las únicas funciones paramétricas de dicha familia de distribuciones que admiten estimadores eficientes son las de la forma:
    \begin{equation*}
        a\cdot \frac{n}{\lm} +b \qquad a,b\in \mathbb{R}, \quad a\neq 0
    \end{equation*}
    y sus estimadores eficientes correspondientes son:
    \begin{equation*}
        a\cdot \sum_{i=1}^{n}X_i + b
    \end{equation*}
    Tomando $a=\nicefrac{1}{n}, b=0$, obtenemos que el estimador $T=\frac{\sum\limits_{i=1}^nX_i}{n}$ es eficiente para $E[X] = \nicefrac{1}{\lm}$, por lo que un Teorema visto en teoría nos asegura que $T$ es el único EMV que podemos considerar para $\nicefrac{1}{\lm}= E[X]$, luego:
    \begin{equation*}
        \wh{E[X]} = \wh{\left(\frac{1}{\lm}\right)} = \frac{\sum\limits_{i=1}^n X_i}{n}
    \end{equation*}

    % // TODO: Qué Teorema?

    Por tanto, por el Teorema de Zehna, tenemos que el EMV de $Var[X]$ es:
    \begin{equation*}
        \wh{Var[X]} = \wh{\left(\frac{1}{\lm^2}\right)} = \wh{\left(\frac{1}{\lm}\right)}^2 = \left(\frac{\sum\limits_{i=1}^n X_i}{n}\right)^2
    \end{equation*}
    
\end{ejercicio}

\begin{ejercicio}
    Sea $X$ una variable aleatoria con función de densidad de la forma
    \begin{equation*}
        f_\theta(x) = \theta x^{\theta -1}, \qquad 0<x<1
    \end{equation*}
    \begin{enumerate}[label=\alph*)]
        \item Calcular un estimador máximo verosímil para $\theta$.

            Maximizamos la función de verosimilitud de $\theta$:
            \begin{equation*}
                L_{x_1,\ldots,x_n}(\theta) = f^n_\theta(x_1, \ldots, x_n) \stackrel{\text{iid.}}{=}\prod_{i=1}^{n}f_\theta(x_i) = \prod_{i=1}^{n} \theta x_i^{\theta-1} = \theta^n \prod_{i=1}^{n} x_i^{\theta -1}
            \end{equation*}

            Puesto que el logaritmo es una función creciente, podemos maximizar la función de verosimilitud aplicando logaritmo, el cual no modifica la abcisa del máximo:
            \begin{equation*}
                \ln L_{x_1,\ldots,x_n}(\theta) = n\ln \theta + (\theta -1)\sum_{i=1}^{n} \ln x_i
            \end{equation*}

            Como buscamos el máximo, derivamos e igualamos a cero para hayar los puntos críticos:
            \begin{equation*}
                \dfrac{\partial \ln L_{x_1,\ldots,x_n}(\theta)}{\partial  \theta} = \frac{n}{\theta} + \sum_{i=1}^{n}\ln x_i = \frac{n + \theta \sum\limits_{k=1}^{n}\ln x_i}{\theta} = 0 \Longleftrightarrow \theta = \frac{-n}{\sum\limits_{k=1}^{n}\ln x_i}
            \end{equation*}

            Sería necesario comprobar que efectivamente es un máximo, pero se entiende que es así. Por tanto, el EMV de $\theta$ (a esperas de comprobar que es un estimador) es:
            \begin{equation*}
                \widehat{\theta}(X_1, \ldots, X_n) = \frac{-n}{\sum\limits_{i=1}^{n} \ln X_i}
            \end{equation*}

            Notemos que $\theta>0$, y que como $0<X_i<1$, tenemos que $\ln X_i < 0$, por lo que $\sum\limits_{i=1}^{n}\ln X_i < 0$, y por tanto $\widehat{\theta} > 0$. Por tanto, tenemos que $\widehat{\theta}$ es un estimador de $\theta$, luego es EMV para $\theta$.
        \item Deducir dicho estimador a partir de los resultados del Ejercicio~\ref{ej:10_rel4}.

            En el Ejercicio~\ref{ej:10_rel4} vimos que la familia de distribuciones en cuestión es regular, con $0<I_X(\theta)<\infty$. Además, vimos que las únicas funciones paramétricas que admiten un estimador eficiente son las de la forma:
            \begin{equation*}
                a\cdot \frac{-n}{\theta} +b\qquad a,b\in \mathbb{R}, \quad a\neq 0
            \end{equation*}
            y sus estimadores eficientes son:
            \begin{equation*}
                a\cdot \sum_{i=1}^{n}\ln X_i + b
            \end{equation*}
            Tomando $a=\nicefrac{-1}{n},b=0$, tenemos que el estimador:
            \begin{equation*}
                T = \frac{\sum\limits_{i=1}^{n}\ln X_i}{-n}
            \end{equation*}
            es eficiente para $\frac{1}{\theta}$, por lo que $T$ es el EMV de $\frac{1}{\theta}$.
            % // TODO: Qué Teorema? Aquí lo define solo para theta
            Si aplicamos el Teorema de Zehna, tenemos que:
            \begin{equation*}
                \widehat{\left(\frac{1}{\theta}\right)} = \frac{1}{\widehat{\theta}} \Longrightarrow \widehat{\theta} = \frac{1}{\widehat{\left(\frac{1}{\theta}\right)} } = \frac{-n}{\sum\limits_{i=1}^{n}\ln X_i}
            \end{equation*}
            Observamos que obtenemos el mismo resultado que en el primer apartado.
    \end{enumerate}
\end{ejercicio}

\begin{ejercicio}
    Sea $(X_1, \ldots, X_n)$ una muestra de una variable $X\rightsquigarrow \cc{B}(k_0,p)$ para cierto $k_0\in \mathbb{N}$ y $p\in \left]0,1\right[$. Estimar, por máxima verosimilitud y por el método de los momentos, el parámetro $p$ y la varianza de $X$.\\

    \noindent
    \textit{Aplicación}: Se lanza $10$ veces un dado cargado y se cuenta el número de veces que sale un $4$. Este experimento se realiza $100$ veces de forma independiente, obteniéndose los siguientes resultados:
    \begin{table}[H]
    \centering
    \begin{tabular}{l|cccc}
        nº de 4 & 0 & 1 & 2 & 3 \\
        \hline
        frecuencia & 84 & 15 & 1 & 0
    \end{tabular}
    \end{table}
    \noindent
    Estimar, a partir de estos datos, la probabilidad de salir un cuatro.\\

    \begin{description}
        \item [Por máxima verosimilitud.] 
            Tenemos:
            \begin{align*}
                P[X_1 =x_1, \ldots, X_n = x_n] &\stackrel{\text{iid.}}{=}\prod_{i=1}^{n} P[X=x_i] = \prod_{i=1}^{n} \binom{k_0}{x_i} p^{x_i}{(1-p)}^{k_0-x_i} \\ &= p^{\sum\limits_{i=1}^{n}x_i}{(1-p)}^{nk_0-\sum\limits_{i=1}^{n}x_i} \prod_{i=1}^{n}\binom{k_0}{n}
            \end{align*}
            
            Puesto que el logaritmo es una función creciente, podemos maximizar la función de verosimilitud aplicando logaritmo, el cual no modifica la abcisa del máximo:
            \begin{equation*}
                \ln L_{x_1,\ldots,x_n}(p) = \ln p \sum_{i=1}^{n}x_i + \ln(1-p)\left(nk_0 - \sum_{i=1}^{n}x_i\right) + \sum_{k=1}^{n}\ln\binom{k_0}{n}
            \end{equation*}
            
            Como buscamos el máximo, derivamos e igualamos a cero para hayar los puntos críticos:
            \begin{align*}
                \dfrac{\partial \ln L_{x_1,\ldots,x_n}(p)}{\partial p} &= \frac{1}{p}\sum_{i=1}^{n}x_i - \frac{1}{1-p}\left(nk_0 - \sum_{i=1}^{n}x_i\right) = \frac{(1-p)\sum\limits_{i=1}^{n}x_i - pnk_0 + p\sum\limits_{i=1}^{n}x_i}{p(1-p)} \\
                                                                       &= \frac{\sum\limits_{i=1}^{n}x_i -pnk_0}{p(1-p)} = 0 \Longleftrightarrow p = \frac{\sum\limits_{i=1}^{n}x_i}{nk_0}
            \end{align*}
            Como $\sum\limits_{i=1}^{n}x_i\leq nk_0$, tenemos entonces que $p\in [0,1]$, por lo que esta fórmula nos dará una estimador, con lo que el EMV es:
            \begin{equation*}
                \widehat{p} = \frac{1}{nk_0}\sum_{i=1}^{n}X_i = \frac{1}{k_0}\overline{X}
            \end{equation*}
            Para $Var[X] = k_0p(1-p)$, si aplicamos el Teorema de Zehna tenemos que:
            \begin{equation*}
                \widehat{Var[X]} = \widehat{k_0p(1-p)} = k_0\widehat{p}(1-\widehat{p}) = \frac{1}{n}\sum_{i=1}^{n}X_i\left(1-\frac{1}{nk_0}\sum_{i=1}^{n}X_i\right) = \overline{X}\left(1-\frac{1}{k_0}\overline{X}\right)
            \end{equation*}
        \item [Por el método de los momentos.] 
            Por el método de los momentos:
            \begin{equation*}
                \left\{\begin{array}{l}
                        k_0p = E[X] = \overline{X} = \frac{1}{n}\sum\limits_{i=1}^{n}X_i \\
                        k_0p(1-p) = Var[X] = A_1 - \overline{X}^2 = \frac{1}{n}\sum\limits_{i=1}^{n}X_i^2 - \overline{X}^2
                \end{array}\right.
            \end{equation*}
            De la primera ecuación obtenemos que:
            \begin{equation*}
                \widehat{p} = \frac{1}{nk_0}\sum_{i=1}^{n}X_i = \frac{1}{k_0}\overline{X}
            \end{equation*}
            De la segunda ecuación obtenemos que:
            \begin{equation*}
                \widehat{Var[X]} = \frac{1}{n}\sum\limits_{i=1}^{n}X_i^2 - \overline{X}^2
            \end{equation*}
    \end{description}
    Para estimar a partir de los datos mencionados la probabilidad de salir un cuatro, lo que hacemos primero es considerar:
    \begin{equation*}
        Y \equiv \text{``Número de veces que sale un 4 en 10 tiradas del dado''} \rightsquigarrow \cc{B}(10,p)
    \end{equation*}
    donde estimaremos el valor de $p$ mediante la estimación obtenida:
    \begin{equation*}
        p= \frac{1}{nk_0}\sum_{i=1}^{n}y_i = \frac{1}{1000}(0\cdot 84 + 1\cdot 15 + 2\cdot 1 + 3\cdot 0) = \frac{15+2}{1000} = \frac{17}{1000} = 0.017
    \end{equation*}
    Por tanto, la estimación de la probabilidad de que salga un cuatro es $0.017$.
\end{ejercicio}

\begin{ejercicio}\label{ej:5_rel5}
    Se lanza un dado hasta que salga un $4$ y se anota el número de lanzamientos necesarios; este experimento se efectúa veinte veces de forma independiente. A partir de los resultados obtenidos, estimar la probabilidad de sacar un $4$ por máxima verosimilitud.\\

    \noindent
    Sea $X$ una variable aleatoria que representa el número de lanzamientos hasta que salga un $4$, es decir:
    \begin{equation*}
        X \equiv \text{``número de lanzamientos antes del lanzamiento en el que sale 4''}
    \end{equation*}
    Tenemos que, si $p$ es la probabilidad de que salga un 4, entonces $X$ sigue una distribución geométrica de parámetro $p$: $X\rightsquigarrow G(p)$. Como se realiza el experimento 20 veces de forma independientes, disponemos de una muestra $x_1, \ldots, x_n$ donde $n=20$. Estimamos $p$ por máxima verosimilitud:
    \begin{equation*}
        L_{x_1,\ldots,x_{n}}(p) = P_p[X_1 =x_1, \ldots, X_n=x_n] \stackrel{\text{indep.}}{=} \prod_{i=1}^{n} P_p[X=x_i] = \prod_{i=1}^{n} {(1-p)}^{x_i}p = p^{n}{(1-p)}^{\sum\limits_{i=1}^{n}x_i}
    \end{equation*}
    
    Como el logaritmo es una función creciente, podemos maximizar la función de verosimilitud aplicando logaritmo, el cual no modifica la abcisa del máximo:
    \begin{equation*}
        \ln L_{x_1,\ldots,x_n}(p) = n\cdot \ln p + \ln(1-p)\sum_{i=1}^{n}x_i
    \end{equation*}
    
    Derivando con respecto a $p$ e igualando a cero para hallar los puntos críticos:
    \begin{equation*}
        \dfrac{\partial \ln L_{x_1,\ldots,x_n}(p)}{\partial p} = \frac{n}{p} - \frac{\sum\limits_{i=1}^{n}x_i}{1-p} = \frac{n - np - p\left(\sum\limits_{i=1}^{n}x_i\right)}{p(1-p)} = \frac{n - p\left(\sum\limits_{i=1}^{n}x_i +n\right)}{p(1-p)}
    \end{equation*}
    Obtenemos por tanto que:
    \begin{equation*}
        \dfrac{\partial \ln L_{x_1,\ldots,x_n}(p)}{\partial  p} = 0 \Longleftrightarrow p = \frac{n}{\sum\limits_{i=1}^{n}x_i + n}
    \end{equation*}
    Como siempre se tiene que dicho valor está en $[0,1]$, tenemos que es un estimador, con lo que el EMV es:
    \begin{equation*}
        \widehat{p} = \frac{n}{\sum\limits_{i=1}^{n}X_i+n} = \frac{20}{\sum\limits_{i=1}^{20}X_i+20} = \frac{1}{1+\ol{X}}
    \end{equation*}
\end{ejercicio}

\begin{ejercicio}
    En $20$ días muy fríos, una granjera pudo arrancar su tractor en el primer, tercer, quinto, primer, segundo, primer, tercer, séptimo, segundo, cuarto, cuarto, octavo, primer, tercer, sexto, quinto, segundo, primer, sexto y segundo intento. Suponiendo que la probabilidad de arrancar en cada intento es constante, y que las observaciones se han obtenido de forma independiente, dar la estimación más verosímil de la probabilidad de que el tractor arranque en el segundo intento.\\

    \noindent
    Reunimos en primer lugar los datos proporcionados por el enunciado en la Tabla~\ref{tab:ej6_rel5}.
    \begin{table}[h]
    \centering
    \begin{tabular}{cccccccc}
        Primer & Segundo & Tercero & Cuarto & Quinto & Sexto & Séptimo & Octavo \\
        \hline
        5 & 4 & 3 & 2 & 2 & 2 & 1 & 1
    \end{tabular}
    \caption{Veces que arrancó en cada intento.}
    \label{tab:ej6_rel5}
    \end{table}

    Consideramos la variable aleatoria dada por:
    \begin{equation*}
        X \equiv \text{``Número de intentos fallidos de arrancar el tractor''}
    \end{equation*}
    Tenemos que $X\rightsquigarrow G(p)$, donde $p$ es la probabilidad de arrancar el tractor en un intento. Tenemos a nuestra disposición de una muestra de $X$, $x_1,\ldots, x_n$, de tamaño $n=20$. Si observamos la similitud con el Ejercicio~\ref{ej:5_rel5}, estamos bajo las mismas hipótesis, por lo que al calcular el EMV para $p$ obtendremos:
    \begin{equation*}
        \widehat{p} = \frac{1}{1+\ol{X}}
    \end{equation*}
    Sin embargo, en este caso no se nos pide estimar $p$, sino la probabilidad de que el tractor arranque en el segundo intento, es decir, $P_p[X=1]$:
    \begin{equation*}
        P_p[X = 1] = (1-p)p
    \end{equation*}
    Si amplicamos el Teorema de Zehna, podemos obtener el EMV de $P_p[X=1]$ a partir de $\widehat{p}$:
    \begin{equation*}
        \widehat{P_p[X=1]} = \widehat{(1-p)p} = \left(1-\widehat{p}\right)\widehat{p}
        = \left(1-\frac{1}{1+\ol{X}}\right)\frac{1}{1+\ol{X}} = \frac{\ol{X}}{1+\ol{X}}\cdot \frac{1}{1+\ol{X}} = \frac{\ol{X}}{(1+\ol{X})^2}
    \end{equation*}

    Por tanto, una vez se ha obtenido el EMV de $P_p[X=1]$, solo nos queda sustituir los valores de la muestra para obtener la estimación. Calculamos en primer lugar la estimación de $\ol{X}$:
    \begin{equation*}
        \ol{X}(x_1, \ldots, x_n) = \frac{1}{n}\sum_{i=1}^{n}x_i = \frac{1}{20}\left(0\cdot 5 + 1\cdot 4 + 2\cdot 3 + 3\cdot 2 + 4\cdot 2 + 5\cdot 2 + 6\cdot 1 + 7\cdot 1\right) = \frac{47}{20} = 2.35
    \end{equation*}

    Por tanto, la estimación máximo verosímil de la probabilidad de que el tractor arranque en el segundo intento es:
    \begin{equation*}
        \widehat{P_p[X=1]}(x_1, \ldots, x_n) = \dfrac{\ol{X}}{(1+\ol{X})^2} = \frac{2.35}{(1+2.35)^2} = \frac{940}{4489} \approx 0.2094
    \end{equation*}
\end{ejercicio}

\begin{ejercicio}
    Una variable aleatoria discreta toma los valores $0$, $1$ y $2$ con las siguientes probabilidades
    \begin{equation*}
        P_p[X=0] = p^2, \qquad P_p[X=1] = 2p(1-p), \qquad P_p[X=2] = {(1-p)}^{2}
    \end{equation*}
    siendo $p$ un parámetro desconocido. En una muestra aleatoria simple de tamaño $100$, se ha presentado $22$ veces el $0$, $53$ veces el $1$ y $25$ veces el $2$. Calcular la función de verosimilitud asociada a dicha muestra y dar la estimación más verosímil de $p$.\\

    Los datos obtenidos son:
    \begin{equation*}
        \begin{array}{c|ccc}
            \text{Dato} & 0 & 1 & 2 \\
            \hline
            \text{Veces} & 22 & 53 & 25
        \end{array}
    \end{equation*}
    
    \begin{description}
        \item[Opción 1.] Sea $X$ una variable aleatoria cuya función masa de probabilidad nos viene dada, tenemos una muestra $X_1, \ldots, X_n$ de tamaño $n=100$. La función de verosimilitud asociada a la muestra es:
        \begin{align*}
            L_{x_1,\ldots,x_n}(p) &= P_p[X_1=x_1, \ldots, X_n = x_n] \stackrel{\text{iid.}}{=}\prod_{i=1}^{n}P_p[X=x_i]\\ 
                                &= \prod_{i=1}^{22}P_p[X=0] \prod_{i=1}^{53}P_p[X=1]\prod_{i=1}^{25}P_p[X=2] \\
                                &= \prod_{i=1}^{22}p^2 \prod_{i=1}^{53}2p(1-p)\prod_{i=1}^{25}{(1-p)}^{2} \\
                                &= p^{44} \cdot 2^{53}p^{53}{(1-p)}^{53} \cdot {(1-p)}^{50} = 2^{53}p^{97}{(1-p)}^{103}
        \end{align*}
        
        Puesto que el logaritmo es una función creciente, podemos maximizar la función de verosimilitud aplicando logaritmo, el cual no modifica la abcisa del máximo:
        \begin{equation*}
            \ln L_{x_1,\ldots,x_n}(p) = 53\ln 2 + 97 \ln p + 103 \ln(1-p)
        \end{equation*}
        
        Derivando con respecto a $p$ e igualando a cero para hallar los puntos críticos:
        \begin{equation*}
            \dfrac{\partial \ln L_{x_1,\ldots,x_n}(p)}{\partial p} = \frac{97}{p} - \frac{103}{1-p} = \frac{97(1-p)-103p}{p(1-p)} = \frac{97 - 200p}{p(1-p)} = 0 \Longleftrightarrow p = \frac{97}{200}
        \end{equation*}

        Por tanto, como $p\in [0,1]$, tenemos que $\frac{97}{200}$ es la estimación máximo verosímil de $p$:
        \begin{equation*}
            \widehat{p} = \frac{97}{200}
        \end{equation*}

        Notemos que este enfoque nos proporciona una estimación máximo verosímil de $p$, pero no nos proporciona el EMV de $p$ como función de la muestra, sino que nos da un valor concreto. Para obtener el EMV de $p$ como función de la muestra, podemos utilizar el enfoque de la Opción 2.
        
        
        \item[Opción 2.] Sea $X$ la variable aleatoria cuya función masa de probabilidad nos viene dada, y vemos que:
        \begin{equation*}
            P[X=k] = \binom{2}{k}(1-p)^kp^{2-k}, \qquad k=0,1,2
        \end{equation*}

        Por tanto, $X\rightsquigarrow B(2,1-p)$, y tenemos una m.a.s. $X_1, \ldots, X_n$ de tamaño $n=100$. La función de verosimilitud asociada a la muestra es:
        \begin{align*}
            L_{x_1,\ldots,x_n}(p) &= P_p[X_1=x_1, \ldots, X_n = x_n] \stackrel{\text{iid.}}{=}\prod_{i=1}^{n}P_p[X=x_i]\\ 
                                &= \prod_{i=1}^{n}\binom{2}{x_i}(1-p)^{x_i}p^{2-x_i} = \prod_{i=1}^{n}\binom{2}{x_i}(1-p)^{\sum\limits_{i=1}^{n}x_i}p^{2n-\sum\limits_{i=1}^{n}x_i}
        \end{align*}

        Puesto que el logaritmo es una función creciente, podemos maximizar la función de verosimilitud aplicando logaritmo, el cual no modifica la abcisa del máximo:
        \begin{equation*}
            \ln L_{x_1,\ldots,x_n}(p) = \sum_{i=1}^{n}\ln\binom{2}{x_i} + \ln(1-p)\sum_{i=1}^{n}x_i + \ln p\left(2n-\sum_{i=1}^{n}x_i\right)
        \end{equation*}
        Derivando con respecto a $p$ e igualando a cero para hallar los puntos críticos:
        \begin{align*}
            \dfrac{\partial \ln L_{x_1,\ldots,x_n}(p)}{\partial p} &= -\frac{\sum\limits_{i=1}^{n}x_i}{1-p} + \frac{2n-\sum\limits_{i=1}^{n}x_i}{p} = \frac{-p\sum\limits_{i=1}^{n}x_i + (1-p)(2n-\sum\limits_{i=1}^{n}x_i)}{p(1-p)}
            =\\&= \frac{2n - 2np - \sum\limits_{i=1}^{n}x_i}{p(1-p)} = 0 \Longleftrightarrow p = \frac{2n-\sum\limits_{i=1}^{n}x_i}{2n} = 1-\frac{\sum\limits_{i=1}^{n}x_i}{2n}
        \end{align*}

        Habría que comprobar que efectivamente es un máximo, pero se entiende que es así. Por tanto, el EMV de $p$ (de nuevo, a esperas de comprobar que es un estimador) es:
        \begin{equation*}
            \widehat{p} = 1-\frac{\sum\limits_{i=1}^{n}X_i}{2n} = 1-\frac{\ol{X}}{2}
        \end{equation*}

        Como $X_i\in \{0,1,2\}$, tenemos que $\ol{X}\in [0,2]$, por lo que $\widehat{p}\in [0,1]$, con lo que $\widehat{p}$ es un estimador de $p$ y por tanto es EMV para $p$. Sustituyendo los valores de la muestra obtenemos:
        \begin{equation*}
            \widehat{p} = 1-\frac{\ol{X}}{2} = 1-\frac{1}{2}\cdot \frac{1}{100}(0\cdot 22 + 1\cdot 53 + 2\cdot 25) = 1-\frac{1}{2}\cdot \frac{103}{100} = 1-\frac{103}{200} = \frac{97}{200}
        \end{equation*}

        Por tanto, la estimación máximo verosímil de $p$ es:
        \begin{equation*}
            \widehat{p} = \frac{97}{200}
        \end{equation*}
    \end{description}
\end{ejercicio}

\begin{ejercicio}
    En el muestreo de una variable aleatoria con distribución $\cc{N}(\mu,1)$, $\mu \in \mathbb{R}$, se observa que no se obtiene un valor menor que $-1$ hasta la quinta observación. Dar una estimación máximo verosímil de $\mu$.\\

    \noindent
    Sea $Y\rightsquigarrow\cc{N}(\mu, 1)$. Consideramos la siguiente variable aleatoria:
    \begin{equation*}
        X\equiv \text{``número de observaciones hasta obtener una menor que $-1$ (sin incluirla)''}
    \end{equation*}

    Considerando que $X$ es el número de ensayos hasta obtener un éxito, siendo el éxito el hecho de que la observación sea menor que $-1$, tenemos que $X\rightsquigarrow G(p)$, donde $p$ es la probabilidad de que una observación sea menor que $-1$, es decir, $p = P[Y<-1]$. Por tanto:
    \begin{equation*}
        P_p[X=x] = {(1-p)}^{x}p
    \end{equation*}

    Calculamos ahora $P[Y<-1]$:
    \begin{align*}
        P[Y<-1] &= P\left[\frac{Y-\mu}{1}<\frac{-1-\mu}{1}\right] = P[Z<-1-\mu]
        =\\&= P[Z>1+\mu] = 1-P[Z<1+\mu] = 1-\Phi(1+\mu)
    \end{align*}
    donde $\Phi$ es la función de distribución de una $\cc{N}(0,1)$. Usando que $\Phi$ es biyectiva, tenemos que:
    \begin{equation*}
        \mu = \Phi^{-1}\left(1-p\right) - 1
    \end{equation*}

    Por tanto, si obtenemos el EMV de $p$, podemos obtener el EMV de $\mu$ aplicando el Teorema de Invarianza de Zehna. Buscamos por tanto el EMV de $p$ a partir de la muestra de $X$. Tenemos que la función de verosimilitud asociada a la muestra es:
    \begin{align*}
        L_{x_1,\ldots,x_n}(p) &= P_p[X_1=x_1, \ldots, X_n=x_n] \stackrel{\text{iid.}}{=}\prod_{i=1}^{n} P_p[X=x_i]
        =\\&= \prod_{i=1}^{n} {(1-p)}^{x_i}p = p^n{(1-p)}^{\sum\limits_{i=1}^{n}x_i}
    \end{align*}

    Como el logaritmo es una función creciente, podemos maximizar la función de verosimilitud aplicando logaritmo, el cual no modifica la abcisa del máximo:
    \begin{equation*}
        \ln L_{x_1,\ldots,x_n}(p) = n\ln p + \ln(1-p)\sum_{i=1}^{n}x_i
    \end{equation*}

    Derivando con respecto a $p$ e igualando a cero para hallar los puntos críticos:
    \begin{equation*}
        \dfrac{\partial \ln L_{x_1,\ldots,x_n}(p)}{\partial p} = \frac{n}{p} - \frac{\sum\limits_{i=1}^{n}x_i}{1-p} = \frac{n - np - p\sum\limits_{i=1}^{n}x_i}{p(1-p)} = \frac{n - p\left(n+\sum\limits_{i=1}^{n}x_i\right)}{p(1-p)} = 0 \Longleftrightarrow p = \frac{n}{n+\sum\limits_{i=1}^{n}x_i}
    \end{equation*}

    Por tanto, de ser un estimador, el EMV de $p$ es:
    \begin{equation*}
        \widehat{p} = \frac{n}{n+\sum\limits_{i=1}^{n}X_i} = \frac{1}{1+\ol{X}}
    \end{equation*}

    Como $X_i\in \mathbb{N}$, tenemos que $\ol{X}\geq 0$, por lo que $\widehat{p}\in [0,1]$, con lo que $\widehat{p}$ es un estimador de $p$ y por tanto es EMV para $p$. Por tanto, aplicando el Teorema de Invarianza de Zehna, tenemos que el EMV de $\mu$ es:
    \begin{equation*}
        \widehat{\mu} = \Phi^{-1}(1-\widehat{p}) - 1
    \end{equation*}

    Evaluando la muestra obtenida, tenemos que $n=1$ y $x_1=5-1=4$, por lo que $\ol{X} = 4$, y por tanto:
    \begin{equation*}
        \widehat{p} = \frac{1}{1+\ol{X}} = \frac{1}{1+4} = \frac{1}{5}
    \end{equation*}

    Por tanto, la estimación máximo verosímil de $\mu$ es:
    \begin{equation*}
        \widehat{\mu} = \Phi^{-1}(1-\widehat{p}) - 1 = \Phi^{-1}\left(1-\frac{1}{5}\right) - 1 = \Phi^{-1}(0.8) - 1
    \end{equation*}

    Si consultamos la tabla de la normal $\cc{N}(0,1)$, obtenemos que $\Phi^{-1}(0.8) \approx 0.84$, por lo que:
    \begin{equation*}
        \widehat{\mu} \approx 0.84 - 1 = -0.16
    \end{equation*}
    es la estimación máximo verosímil de $\mu$.
\end{ejercicio}

\begin{ejercicio}
    En la producción de filamentos eléctricos la medida de interés, $X$, es el tiempo de vida de cada filamento, que tiene una distribución exponencial de parámetro $\theta$. Se eligen $n$ de tales filamentos de forma aleatoria e independiente, pero, por razones de economía, no conviene esperar a que todos se quemen y la observación acaba en el tiempo $T$. Dar el estimador máximo verosímil para la media de $X$ a partir del número de filamentos quemados durante el tiempo de observación.\\


    Sea $X\rightsquigarrow \exp(\theta)$, y consideramos la siguiente variable aleatoria:
    \begin{center}
        $Y\equiv$ ``número de filamentos con tiempo de vida menor que $T$ en una muestra de tamaño $n$''
    \end{center}

    Tenemos que $Y\rightsquigarrow B(n,p)$, donde $p$ es la probabilidad de que el tiempo de vida de un filamento sea menor que $T$, es decir, $p = P[X<T]$. Sabemos que el parámetro a estimar es:
    \begin{equation*}
        E[X] = \frac{1}{\theta}
    \end{equation*}

    Por tanto, por el Teorema de Invarianza de Zehna, si obtenemos el EMV de $\theta$, podemos obtener el EMV de $E[X]$ a partir de él. Calculamos ahora el valor de $p$:
    \begin{align*}
        p &= P[X<T] = \int_{0}^{T}\theta e^{-\theta x}dx = \left[-e^{-\theta x}\right]_{0}^{T} = 1 - e^{-\theta T}
        \Longrightarrow
        e^{-\theta T} = 1-p \Longrightarrow \theta = \frac{-\ln(1-p)}{T}
    \end{align*}

    Por tanto, como el valor de $T$ es conocido, si obtenemos el EMV de $p$, podemos obtener el EMV de $\theta$ a partir de él, y por tanto el EMV de $E[X]$. Buscamos por tanto el EMV de $p$ a partir de la muestra de $Y$, que es de tamaño $1$ (notemos que la muestra de $Y$ es de tamaño $1$ porque solo se realiza un experimento, que consiste en observar el número de filamentos quemados en una muestra de tamaño $n$. No se debe confundir con la muestra de tamaño $n$ de filamentos, que es la que se utiliza para definir $Y$). Por tanto, tenemos que la función de verosimilitud asociada a la muestra es:
    \begin{equation*}
        L_{y_1}(p) = P_p[Y=y_1] = \binom{n}{y}p^{y_1}{(1-p)}^{n-{y_1}}
    \end{equation*}

    Puesto que el logaritmo es una función creciente, podemos maximizar la función de verosimilitud aplicando logaritmo, el cual no modifica la abcisa del máximo:
    \begin{equation*}
        \ln L_{y_1}(p) = \ln\binom{n}{y_1} + y_1\ln p + (n-y_1)\ln (1-p)
    \end{equation*}

    Derivando con respecto a $p$ e igualando a cero para hallar los puntos críticos:
    \begin{equation*}
        \dfrac{\partial \ln L_{y_1}(p)}{\partial p} = \frac{y_1}{p} - \frac{n-y_1}{1-p} = \frac{y_1 - y_1p - np + y_1p}{p(1-p)} = \frac{y_1-np}{p(1-p)} = 0 \Longleftrightarrow p = \frac{y_1}{n}
    \end{equation*}

    Por tanto, de ser un estimador, el EMV de $p$ es:
    \begin{equation*}
        \widehat{p} = \frac{Y}{n}
    \end{equation*}

    Como $Y\in \{0,1,\ldots,n\}$, tenemos que $\widehat{p}\in [0,1]$, con lo que $\widehat{p}$ es un estimador de $p$ y por tanto es EMV para $p$. Por tanto, aplicando el Teorema de Invarianza de Zehna, tenemos que el EMV de $\theta$ es:
    \begin{equation*}
        \widehat{\theta} = \frac{-\ln(1-\widehat{p})}{T} = \frac{-\ln\left(1-\nicefrac{Y}{n}\right)}{T}
    \end{equation*}

    Aplicando de nuevo el Teorema de Invarianza de Zehna, tenemos que el EMV de $E[X]$ es:
    \begin{equation*}
        \widehat{E[X]} = \frac{1}{\widehat{\theta}} = \frac{-T}{\ln(1-\widehat{p})} = \frac{-T}{\ln\left(1-\nicefrac{Y}{n}\right)}
    \end{equation*}
\end{ejercicio}

\begin{ejercicio}
    Sean $X_1, \ldots, X_n$ observaciones independientes de una variable $X\rightsquigarrow\{\Gamma(p,a) : p,a>0\}$. Estimar ambos parámetros mediante el método de los momentos.\\

    \noindent
    \textit{Aplicación}: Ciertos neumáticos radiales tuvieron vidas útiles de $35200$, $41000$, $44700$, $38600$ y $41500$ kilómetros. Suponiendo que estos datos son observaciones independientes de una variable con distribución exponencial de parámetro $\theta$, dar una estimación de dicho parámetro por el método de los momentos.\\

    Planteamos el sistema a resolver, sabiendo que si $X\rightsquigarrow \Gamma(p,a)$, entonces:
    \begin{equation*}
        E[X] = \frac{p}{a}, \qquad Var[X] = \frac{p}{a^2}
    \end{equation*}

    Como tratamos estimar dos parámetros, necesitamos dos ecuaciones. Respecto al momento de primer orden no centrado, tenemos que:
    \begin{equation*}
        E[X] = \frac{p}{a} =A_1 = \ol{X}
    \end{equation*}

    Respecto al momento de segundo orden centrado, tenemos que:
    \begin{equation*}
        Var[X] = \frac{p}{a^2} = A_2 - \ol{X}^2
    \end{equation*}

    Por tanto, tenemos el siguiente sistema de ecuaciones con dos incógnitas:
    \begin{equation*}
        \left\{\begin{array}{rcl}
                \nicefrac{p}{a} &=& \overline{X} \\
                \nicefrac{p}{a^2} &=& A_2 - {\left(\overline{X}\right)}^2
        \end{array}\right.
    \end{equation*}
    
    De la primera ecuación tenemos que:
    \begin{equation*}
        p = a\overline{X}
    \end{equation*}
    
    Si sustituimos en la segunda, obtenemos el valor de $a$:
    \begin{equation*}
        \frac{\overline{X}}{a} = A_2 - {\left(\overline{X}\right)}^2 \Longrightarrow a = \frac{\overline{X}}{A_2-{\left(\overline{X}\right)}^2}
    \end{equation*}

    Sustituyendo ahora $a$ en la expresión de $p$, obtenemos el valor de $p$:
    \begin{equation*}
        p = \frac{{\left(\overline{X}\right)}^{2}}{A_2-{\left(\overline{X}\right)}^{2}}
    \end{equation*}

    Por tanto, los estimadores de $a$ y $p$ por el método de los momentos son:
    \begin{equation*}
        \widehat{a} = \frac{\overline{X}}{A_2-{\left(\overline{X}\right)}^2}, \qquad \widehat{p} = \frac{{\left(\overline{X}\right)}^{2}}{A_2-{\left(\overline{X}\right)}^{2}}
    \end{equation*}~\\

    \noindent
    Sea $X\rightsquigarrow \exp(\theta)$. Notemos que, aunque la exponencial sea un caso particular de la gamma, con $p=1$ y $a=\theta$, no podemos aplicar directamente los estimadores obtenidos para la gamma, ya que en este caso solo tenemos un parámetro a estimar. Por tanto, aplicamos el método de los momentos directamente a la exponencial. Sabemos que:
    \begin{equation*}
        E[X] = \frac{1}{\theta} = A_1 = \ol{X}
    \end{equation*}

    De esta ecuación, despejamos $\theta$:
    \begin{equation*}
        \theta = \frac{1}{\overline{X}}
    \end{equation*}

    Por tanto, la estimación de $\theta$ por el método de los momentos es:
    \begin{equation*}
        \widehat{\theta} = \frac{1}{\overline{X}}
    \end{equation*}

    Calculamos ahora $\overline{X}$:
    \begin{align*}
        \overline{X} &= \frac{35200 + 41000 + 44700 + 38600 + 41500}{5} = 40200
    \end{align*}

    Por tanto, la estimación de $\theta$ por el método de los momentos es:
    \begin{equation*}
        \widehat{\theta} = \frac{1}{\overline{X}} = \frac{1}{40200} \approx 2.488\cdot 10^{-5}
    \end{equation*}






    \begin{comment}
    Notemos que una exponencial es un caso particular de una gamma, con $p=1$ y $a=\theta$. Por tanto, empleando el estimador de $a$ obtenido anteriormente, tenemos que:
    \begin{equation*}
        \widehat{\theta} = \frac{\overline{X}}{A_2-{\left(\overline{X}\right)}^2}
    \end{equation*}

    Calculamos ahora $A_2$ y $\overline{X}$:
    \begin{align*}
        A_2 &= \frac{35200^2 + 41000^2 + 44700^2 + 38600^2 + 41500^2}{5} = 1.626068\cdot 10^9\\
        \overline{X} &= \frac{35200 + 41000 + 44700 + 38600 + 41500}{5} = 40200
    \end{align*}

    Por tanto, la estimación de $\theta$ por el método de los momentos es:
    \begin{equation*}
        \widehat{\theta} = \frac{40200}{1.626068\cdot 10^9 - 40200^2} \approx 4.0088\cdot 10^{-3}
    \end{equation*}
    \end{comment}
\end{ejercicio}


\subsection*{Ejercicios Adicionales}
\begin{ejercicio}[Reto]
    Sea $X_1, \ldots, X_n$ una muestra aleatoria simple de una variable $X$ con función de densidad:
    \begin{equation*}
        f_\theta(x) = \dfrac{-2x}{(1-\theta)^2}\qquad 1-\theta < x < 0
    \end{equation*}

    Calcular el EMV de $\theta$ por el método de máxima verosimilitud y por el método de los momentos.
    \begin{description}
        \item[Por el Método de los Momentos]~
        
        Buscamos el primer momento de la variable $X$:
        \begin{align*}
            E[X] &= \int_{1-\theta}^{0} x f_\theta(x) dx = \int_{1-\theta}^{0} x \frac{-2x}{(1-\theta)^2} dx = \frac{-2}{(1-\theta)^2}\int_{1-\theta}^{0} x^2 dx = \frac{-2}{(1-\theta)^2}\left[\frac{x^3}{3}\right]_{1-\theta}^{0} =\\&= \frac{-2}{(1-\theta)^2}\left(0 - \frac{(1-\theta)^3}{3}\right) = \frac{2(1-\theta)}{3}
        \end{align*}

        Por el método de los momentos, igualamos la media muestral con la media poblacional:
        \begin{equation*}
            \overline{X} = E[X] \Longrightarrow \overline{X} = \frac{2(1-\theta)}{3} \Longrightarrow 1-\theta = \frac{3\overline{X}}{2} \Longrightarrow \widehat{\theta} = 1 - \frac{3\overline{X}}{2}
        \end{equation*}

        Como $\theta>1$ y $\overline{X}<0$, tenemos que $\widehat{\theta}>1$, luego la estimación obtenida es un estimador.

        
        \item[Por el Método de Máxima Verosimilitud]~
        
        La función de verosimilitud es:
        \begin{equation*}
            L_{x_1,\ldots,x_n}(\theta) = \prod_{i=1}^{n} f_\theta(x_i) = \prod_{i=1}^{n} \frac{-2x_i}{(1-\theta)^2} = \frac{(-2)^n \prod_{i=1}^{n} x_i}{(1-\theta)^{2n}}\qquad \theta < 1 - x_{(1)}
        \end{equation*}
        \begin{itemize}
            \item Si $\theta \leq 1 + x_{(1)}$, entonces $L_{x_1,\ldots,x_n}(\theta) = 0$.
            \item Si $\theta > 1 + x_{(1)}$, entonces $L_{x_1,\ldots,x_n}(\theta) > 0$.
        \end{itemize}

        Veamos que la función de verosimilitud es creciente en el intervalo $\theta > 1 - x_{(1)}$:
        \begin{equation*}
            \frac{\partial L_{x_1,\ldots,x_n}(\theta)}
            {\partial \theta} = \frac{(-2)^n \prod\limits_{i=1}^{n} x_i}{(1-\theta)^{2n+1}} \cdot 2n >0
        \end{equation*}
        Por tanto, el EMV de $\theta$ es:
        \begin{equation*}
            \widehat{\theta} = 1 - X_{(1)}
        \end{equation*}

        Como $\theta>1$ y $X_{(1)}<0$, tenemos que $\widehat{\theta}>1$, luego la estimación obtenida es un estimador.
    \end{description}
\end{ejercicio}



\begin{ejercicio}
    Sea $X_1, \ldots, X_n$ una muestra aleatoria simple de una variable aleatoria $X$, que sigue una distribución de Rayleigh de parámetro $\theta > 0$, con función de densidad dada por:
    \begin{equation*}
        f_\theta(x) = \frac{x}{\theta^2} \exp\left(-\frac{x^2}{2\theta^2}\right), \qquad x \geq 0
    \end{equation*}
    \begin{enumerate}[label=\alph*)]
        \item Calcular el estimador de máxima verosimilitud de $\theta$.
        
        Calculamos en primer lugar la función de verosimilitud para la muestra $(x_1, \ldots, x_n)$:
        \begin{align*}
            L_{x_1, \ldots, x_n}(\theta) &= \prod_{i=1}^{n} f_\theta(x_i) = \prod_{i=1}^{n} \frac{x_i}{\theta^2} \exp\left(-\frac{x_i^2}{2\theta^2}\right) = \\
            &= \left(\prod_{i=1}^{n} x_i\right) \dfrac{1}{\theta^{2n}} \exp\left(-\frac{1}{2\theta^2}\sum_{i=1}^{n} x_i^2\right) \qquad \forall \theta > 0.
        \end{align*}

        Puesto que el logaritmo es una función estrictamente creciente, podemos trabajar con la log-verosimilitud:
        \begin{equation*}
            \ln L_{x_1, \ldots, x_n}(\theta) = \sum_{i=1}^{n} \ln x_i - 2n \ln \theta - \frac{1}{2\theta^2} \sum_{i=1}^{n} x_i^2 \qquad \forall \theta > 0.
        \end{equation*}

        Como buscamos el máximo, derivamos respecto a $\theta$ e igualamos a cero:
        \begin{align*}
            \frac{d}{d\theta}\ln L_{x_1, \ldots, x_n}(\theta) &= -\frac{2n}{\theta} + \frac{1}{\theta^3}\sum_{i=1}^{n} x_i^2 = 0 \implies \\
            &\implies -2n\theta^2 + \sum_{i=1}^{n} x_i^2 = 0 \implies \theta^2 = \frac{1}{2n}\sum_{i=1}^{n} x_i^2.
        \end{align*}

        Dado que $\theta > 0$, despejamos:
        \begin{equation*}
            \theta = \sqrt{\frac{1}{2n}\sum_{i=1}^{n} x_i^2}.
        \end{equation*}

        Sería necesario comprobar que efectivamente se trata de un máximo, pero se tiene. Por tanto, el candidato a EMV de $\theta$ es:
        \begin{equation*}
            \widehat{\theta} = \sqrt{\frac{1}{2n}\sum_{i=1}^{n} X_i^2}.
        \end{equation*}

        Como $\theta > 0$ y $\widehat{\theta} > 0$, podemos concluir que efectivamente $\widehat{\theta}$ es un estimador de $\theta$, y por lo tanto, $\widehat{\theta}$ es el EMV de $\theta$.
        \item Dada una constante conocida $c > 0$, calcular el estimador de máxima verosimilitud de $P_\theta[X > c]$.
        
        En primer lugar, calculamos la probabilidad teórica para $c > 0$:
        \begin{align*}
            P_\theta[X > c] &= \int_{c}^{\infty} \frac{x}{\theta^2} \exp\left(-\frac{x^2}{2\theta^2}\right) dx = \left[ -\exp\left(-\frac{x^2}{2\theta^2}\right) \right]_{c}^{\infty} = \\
            &= 0 - \left(-\exp\left(-\frac{c^2}{2\theta^2}\right)\right) = \exp\left(-\frac{c^2}{2\theta^2}\right).
        \end{align*}

        Por tanto, buscamos el EMV de la función paramétrica:
        \begin{equation*}
            g(\theta) = P_\theta[X > c] = \exp\left(-\frac{c^2}{2\theta^2}\right).
        \end{equation*}

        Por el Teorema de Invarianza de Zehna, si $\widehat{\theta}$ es el EMV de $\theta$, entonces:
        \begin{equation*}
            g(\widehat{\theta}) = \exp\left(-\frac{c^2}{2\widehat{\theta}^2}\right)
        \end{equation*}
        es el EMV de $g(\theta) = P_\theta[X > c]$.

        Sustituyendo la expresión obtenida para $\widehat{\theta}$ en el apartado anterior:
        \begin{equation*}
            \widehat{\theta}^2 = \frac{1}{2n}\sum_{i=1}^{n} X_i^2,
        \end{equation*}
        obtenemos que el EMV de $P_\theta[X > c]$ es:
        \begin{equation*}
            \wh{P_\theta[X > c]} = \exp\left(-\frac{c^2}{2\left(\frac{1}{2n}\sum\limits_{i=1}^{n} X_i^2\right)}\right) = \exp\left(-\frac{n c^2}{\sum\limits_{i=1}^{n} X_i^2}\right).
        \end{equation*}
    \end{enumerate}

\end{ejercicio}


\begin{ejercicio}
    Sea $(X_1, \dots, X_n)$ una muestra aleatoria simple de una variable aleatoria $X \rightsquigarrow \mathcal{N}(\mu, \mu^2)$ con $\mu > 0$.
    \begin{enumerate}
        \item Calcula los estimadores de $\mu$ basados en el primer y en el segundo momento muestral por separado usando el método de los momentos.
        
        Sabemos que $E[X] = \mu$, $Var[X] = \mu^2$ y
        \begin{equation*}
            m_2 = E[X^2] = Var[X] + (E[X])^2 = \mu^2 + \mu^2 = 2\mu^2
        \end{equation*}

        Si usamos el \textbf{primer momento}:
        \begin{equation*}
            m_1 = A_1 \implies \mu = \overline{X} \implies \widehat{\mu}_1 = \overline{X}
        \end{equation*}

        Si usamos el \textbf{segundo momento}:
        \begin{equation*}
            m_2 = A_2 \implies 2\mu^2 = \frac{1}{n}\sum_{i=1}^n X_i^2 \implies \widehat{\mu}_2 = \sqrt{\frac{1}{2n}\sum_{i=1}^n X_i^2}
        \end{equation*}

        Notemos que obtenemos dos estimadores distintos de $\mu$ a partir de los dos momentos muestrales.


        \item Plantea la ecuación de máxima verosimilitud para $\mu$ y deduce el EMV. Demuestra algebraicamente cómo este EMV ``fusiona'' la información de los dos estimadores del método de los momentos que hallaste en el primer apartado.
        
        La función de verosimilitud es:
        \begin{equation*}
            L_{x_1,\hdots,x_n}(\mu) \overset{\text{iid}}{=} \prod_{i=1}^n \frac{1}{\sqrt{2\pi\mu^2}} \exp\left( -\frac{(x_i - \mu)^2}{2\mu^2} \right) = (2\pi\mu^2)^{\nicefrac{-n}{2}} \exp\left( -\frac{1}{2\mu^2} \sum_{i=1}^n (x_i - \mu)^2 \right)
        \end{equation*}

        Desarrollamos el cuadrado en el sumatorio: 
        \begin{equation*}
            \sum_{i=1}^n (x_i - \mu)^2 = \sum_{i=1}^n x_i^2 - 2\mu \sum_{i=1}^n x_i + n\mu^2
        \end{equation*}

        Puesto que el logaritmo es una función estrictamente creciente, podemos trabajar con la log-verosimilitud:
        \begin{equation*}
            \ln L(\mu) = -\frac{n}{2}\ln(2\pi) - n\ln\mu - \frac{1}{2\mu^2}\sum_{i=1}^n x_i^2 + \frac{1}{\mu}\sum_{i=1}^n x_i - \frac{n}{2}
        \end{equation*}

        Derivamos respecto a $\mu$ e igualamos a cero:
        \begin{equation*}
            \frac{\partial \ln L(\mu)}{\partial\mu} = -\frac{n}{\mu} + \frac{1}{\mu^3}\sum_{i=1}^n x_i^2 - \frac{1}{\mu^2}\sum_{i=1}^n x_i = 0
        \end{equation*}

        Multiplicando por $\mu^3$ y dividiendo entre $n$, obtenemos:
        \begin{equation*}
            -\mu^2 + \frac{1}{n}\sum_{i=1}^n x_i^2 - \mu\frac{1}{n}\sum_{i=1}^n x_i = 0 \Longleftrightarrow \mu^2 + \overline{X}\mu - A_2 = 0
        \end{equation*}

        Despejando $\mu$ de la ecuación cuadrática, obtenemos:\begin{equation*}
            \mu = \frac{-\overline{X} \pm \sqrt{\overline{X}^2 + 4A_2}}{2}
        \end{equation*}

        Como $A_2 \ge 0$, el discriminante siempre es mayor que $\overline{X}^2$, por lo que la raíz cuadrada es mayor que $|\overline{X}|$. Como el parámetro $\mu$ está definido estrictamente positivo ($\mu > 0$), la única solución matemáticamente posible es tomar el signo positivo en la ecuación anterior, por lo que el EMV de $\mu$ es:
        \begin{equation*}
            \widehat{\mu} = \frac{-\overline{X} + \sqrt{\overline{X}^2 + 4A_2}}{2}
        \end{equation*}

        Notemos además que se trata de un estimador, ya que $\widehat{\mu} > 0$ para cualquier muestra $(X_1, \dots, X_n)$. Si recordamos del primer apartado que $\widehat{\mu}_1 = \overline{X}$ y $A_2 = 2\widehat{\mu}_2^2$, el EMV se puede reescribir como:
        \begin{equation*}
            \widehat{\mu} = \frac{-\widehat{\mu}_1 + \sqrt{\widehat{\mu}_1^2 + 8\widehat{\mu}_2^2}}{2}
        \end{equation*}

        El EMV de $\mu$ es una combinación de los dos estimadores obtenidos por el método de los momentos, lo que significa que el EMV ``fusiona'' la información de ambos estimadores. En particular, el EMV toma en cuenta tanto la media muestral como el segundo momento muestral para proporcionar una estimación más precisa del parámetro $\mu$.


        \item Aplica el Teorema de Invarianza de Zehna para encontrar el EMV de la probabilidad de que una medición aleatoria sea negativa, es decir, estimar $P_\mu[X < 0]$. Interpreta el resultado.
        
        La probabilidad de que una medición aleatoria sea negativa es:
        \begin{equation*}
            P_\mu[X < 0] = \Phi\left(\frac{0 - \mu}{\mu}\right) = \Phi(-1) \approx 0.1587
        \end{equation*}
        donde $\Phi$ es la función de distribución acumulada de la normal estándar. Notemos que esta probabilidad no depende de $\mu$, sino que es una constante.

        Por el Teorema de Invarianza de Zehna, si $\widehat{\mu}$ es el EMV de $\mu$, entonces:
        \begin{equation*}
            \widehat{P_\mu[X < 0]} = P_{\widehat{\mu}}[X < 0] = \Phi\left(\frac{0 - \widehat{\mu}}{\widehat{\mu}}\right) = \Phi(-1) \approx 0.1587
        \end{equation*}

        El estimador máximo verosímil de esta probabilidad no depende en absoluto de la muestra $(X_1, \dots, X_n)$. La propia estructura del modelo $\mathcal{N}(\mu, \mu^2)$ impone matemáticamente que siempre habrá un $\approx 15.866\%$ de masa de probabilidad en el eje negativo, independientemente de la magnitud de la media.
    \end{enumerate}
\end{ejercicio}